\section{The product cut and bounded-degree descent}
\label{sec:product-cut}

The positive index from Proposition~\ref{prop:index} will produce a
hypersurface containing some $y_{i\ell}$ but not $Z_{i\ell}$.
Both its degree and
the degree of a field of definition must be bounded: geometric degree
alone would not contradict Proposition~\ref{prop:minimal-products}.
We first transfer the section to a polynomial on a product of
projective spaces, and then keep track of all conjugates of a suitable
component of its index locus.

\subsection{Finite projections and the norm of a section}

The section $s_i$ supplied by Lemma~\ref{lemma:small-sections} is
defined over $k_i$.  The larger field used in the local index proof
serves only to evaluate this section and its coefficients at the
marked point; it does not change the section's field of definition.

Choose a fixed positive integer $t_0$ such that
\begin{equation}\label{eq:cut-ample-difference}
 Q_i=t_0\sum_\ell b_{i\ell}^2L_\ell-F_i^*L
\end{equation}
is ample on $\cA$.  This is possible uniformly in $i$: after changing
variables by $\operatorname{diag}(b_{i\ell})$ in the invariant
Hermitian forms, the first term is a fixed positive form times $t_0$,
whereas the second is bounded because
$F_i\operatorname{diag}(b_{i\ell}^{-1})$ is bounded.

Take a positive power making $Q_i$ globally generated.  Replacing
$s_i$ by the corresponding power, and $n_i$ by the same multiple,
choose a section of $Q_i^{n_i}$ over $k_i$ nonzero at $y_i$ and
multiply.  Such a section exists over $k_i$: among a $k_i$-basis of
a generating space, some member has nonzero value at the algebraic
point $y_i$.  We obtain a nonzero section $s'_i$ of
\[
 \left.\boxtimes_\ell L_\ell^{t_0n_ib_{i\ell}^2}\right|_{Z_i}
\]
on the resolution.  Its index with weights $n_ib_{i\ell}^2$ is still
at least the fixed constant of Proposition~\ref{prop:index}; taking
a power scales numerator and denominator together, and multiplication
by a regular section cannot lower weighted order.

Let $Z_i^\nu$ be the normalization of $Z_i$.  The proper
birational map from the resolution to $Z_i^\nu$ has
pushforward structure sheaf $\mathcal O_{Z_i^\nu}$.  The
projection formula therefore descends $s'_i$ to the pullback of the
displayed bundle on $Z_i^\nu$.  We do not require descent to
the possibly nonnormal variety $Z_i$ itself.

For each factor choose a finite linear projection
\[
 \pi_{i\ell}:Z_{i\ell}\longrightarrow\mathbb P^{d_\ell},
 \qquad \pi_{i\ell}^*\mathcal O(1)=L_\ell|_{Z_{i\ell}},
\]
which is \'{e}tale at $y_{i\ell}$.  A general projection has this
property at a specified smooth point.  The same requirements at its
finitely many conjugates give a nonempty open condition defined over
$k_i$ in the space of projections.  Since $k_i$ is infinite, it
contains a choice over $k_i$.  The degree of this finite map is
$\deg_{L_\ell}Z_{i\ell}$ and is bounded.  A zero-dimensional
geometrically integral factor defined over $k_i$ is a $k_i$-point;
its projection to $\mathbb P^0$ has degree one.

The product projection, composed with normalization, is a finite map
\[
 \Pi_i:Z_i^\nu\longrightarrow
            P_i=\prod_\ell\mathbb P^{d_\ell}
\]
of bounded degree
$e_i=\prod_\ell\deg_{L_\ell}Z_{i\ell}$.
Take the field norm of $s'_i$ along this map, as in
\cite[p.~565]{Faltings1991}.  It is a nonzero
section $G_i$ on $P_i$, defined over $k_i$, with multidegrees
\begin{equation}\label{eq:cut-multidegrees}
 D_{i\ell}=e_it_0n_ib_{i\ell}^2.
\end{equation}
For completeness, on an affine trivializing open of the target,
$s'_i$ is an element of a finite algebra over the coordinate ring.
It is integral there; the coefficients of its characteristic
polynomial in the function-field extension are integral and belong
to the fraction field of the normal target ring, and hence belong
to that ring.  Thus its norm is regular.  On overlaps, a change of
frame is raised to the power $e_i$, proving
\eqref{eq:cut-multidegrees}.  This also proves that the construction
works without a flatness hypothesis for $\Pi_i$.

The characteristic polynomial shows that the pullback of $G_i$ is
divisible by $s'_i$ in the local ring at $y_i$.  Because the
projections are \'{e}tale factor by factor there, their completed local
rings over $\Qbar$ identify through factorwise parameters.  If
$z_i=\Pi_i(y_i)$, it follows that
\begin{equation}\label{eq:cut-polynomial-index}
 \operatorname{ind}_{D_i}(G_i;z_i)\ge\eta_1>0,
\end{equation}
where $\operatorname{ind}_{D_i}$ gives parameters in the $\ell$th
projective factor weight $D_{i\ell}^{-1}$.  Indeed the change from
the old weights divides the index by $e_it_0$, which is bounded.
We now omit the zero-dimensional projective factors.  At least one
factor remains because $d_0>0$.  For the remaining ordered factors,
\begin{equation}\label{eq:cut-ratios}
 D_{i\ell}/D_{i,\ell+1}\longrightarrow\infty,
\end{equation}
including when omitted factors occur between them.

\subsection{Components of separated index loci}

Let $G$ be a nonzero multihomogeneous polynomial of positive
multidegrees $D=(D_1,\ldots,D_m)$ on
$P=\prod_{\ell=1}^m\mathbb P^{d_\ell}$, with $d_\ell\ge1$, and write
$N=\sum_\ell d_\ell$.  Its index at a point is the minimum of
$\sum_\ell|\alpha_\ell|/D_\ell$ among nonzero coefficients in
local affine coordinates and a local frame.  Equivalently one may
use homogeneous partial derivatives, with their total orders in
each block as weights.  This equivalence follows by dehomogenizing
in a nonzero coordinate: affine derivatives are homogeneous
derivatives divided by the local frame, while derivatives in the
omitted coordinate are combinations of affine derivatives of no
larger weighted order, by homogeneity.

For a positive irrational number $\sigma$, let $\mathcal I_\sigma$
be the closed locus where this index exceeds $\sigma$.  It is cut
out by all homogeneous derivatives having weighted order less than
$\sigma$.  All occurring weights are rational, so changing either
strict inequality to a nonstrict one gives the same locus.

Suppose that $z\in P$ has index at least $\eta>0$.  Following
\cite[Remark~3.4]{Faltings1991}, we first find a component through $z$
that persists between two separated levels.
Choose a positive irrational $\sigma_0$ and a positive rational
$\epsilon\le1$ such that the $N+2$ levels
\[
 \sigma_j=\sigma_0+j\epsilon\qquad(0\le j\le N+1)
\]
all lie below $\eta$.  These choices depend only on $\eta$ and $N$,
not on the multidegrees $D$.  The point $z$ lies in every corresponding
index locus.  Choose a component $J_{N+1}$ through $z$ at the highest
level and then, at each lower level, a component containing the one
already chosen.  This gives
\[
 J_{N+1}\subseteq J_N\subseteq\cdots\subseteq J_0,
 \qquad J_j\text{ a component of }\mathcal I_{\sigma_j}.
\]
Every $J_j$ is proper in $P$, because $G\ne0$ and all levels are
positive.  Two consecutive components coincide: otherwise their
dimensions would strictly increase more than $N$ times.  Their common
value $J$ contains $z$ and is a component of both
$\mathcal I_\sigma$ and $\mathcal I_{\sigma+\epsilon}$ for one of
the chosen levels $\sigma$.

We use Faltings' product theorem \cite[Theorem~3.1]{Faltings1991}
in Evertse's explicit form:
for $m\ge2$, $0<\epsilon\le1$, and
\begin{equation}\label{eq:cut-product-threshold}
 D_\ell/D_{\ell+1}\ge(mN/\epsilon)^N,
\end{equation}
a common irreducible component of the index loci at levels
$\sigma$ and $\sigma+\epsilon$ is a product of subvarieties of the
individual projective factors.  This is
\cite[Theorem~1]{Evertse1995Product}, an explicit version of
Faltings' theorem; its hypotheses concern a nonzero polynomial over
an algebraically closed field of characteristic zero.  Thus, over
$\Qbar$, the component just constructed is a product whenever
\eqref{eq:cut-product-threshold} holds.  For $m=1$ it is already a
subvariety of the single projective factor.

It remains to bound both the factor degrees and a common field of
definition independently of $D$.  Evertse's theorem also gives such
an arithmetic bound.  The next lemma proves the needed weaker bound
directly by Faltings' conjugate-counting method
\cite[Remark~3.2]{Faltings1991}, including the case $m=1$.

\begin{lemma}\label{lemma:conjugate-product-degree}
Suppose that $G$ is defined over a number field $k$, and that a
proper geometric component
$J=\prod_{\ell=1}^mJ_\ell$ is common to
$\mathcal I_\sigma$ and $\mathcal I_{\sigma+\epsilon}$, where
$\sigma>0$ and $0<\epsilon\le1$, and neither level is a rational
derivative weight.  If $q$ is the number of conjugates of $J$ over
$k$, then
\begin{equation}\label{eq:conjugate-product-degree}
 q\prod_\ell\deg J_\ell\le C(N,\epsilon).
\end{equation}
In particular the factors have bounded degrees and are defined over
a common extension of $k$ of bounded degree.
\end{lemma}

\begin{proof}
Put $c_\ell=\codim_{\mathbb P^{d_\ell}}J_\ell$ and
$c=\sum_\ell c_\ell>0$.  Both index loci are defined over $k$, so
every conjugate of $J$ is again a common component.  All conjugates
have the same factor codimensions.  Choose general linear spaces
$\Lambda_\ell\simeq\mathbb P^{c_\ell}$ and put
$\Lambda=\prod_\ell\Lambda_\ell$.  Each conjugate of $J$ meets
$\Lambda$ transversely in $\prod_\ell\deg J_\ell$ reduced points.
We may choose the slices so that all these sets are disjoint and
none of their points belongs to any other component of
$\mathcal I_\sigma$.

Here is the incidence justification for this simultaneous choice.
The intersection of a component $J'$ with another component of
$\mathcal I_\sigma$ is a proper closed subset of $J'$, hence has
dimension less than $\dim J'$.  For a fixed point, the condition
that it belong to $\Lambda$ has codimension
$\sum_\ell(d_\ell-c_\ell)=\dim J'$ in the product of the
Grassmannians parametrizing the slices.  Thus a general product
slice misses every such proper subset.  There are only finitely
many components and conjugates.  The same argument, together with
generic transversality in each factor in characteristic zero,
avoids the singular loci.  It follows that the indicated
$q\prod_\ell\deg J_\ell$ points are isolated in
$\Lambda\cap\mathcal I_\sigma$.  They all belong to the higher
index locus.

Multiply each derivative defining $\mathcal I_\sigma$ by all
monomials of the missing multidegree, so that every resulting
equation has multidegree $D$.  This operation does not change their
common zero set: at any point, some monomial of each specified
nonnegative multidegree is nonzero.  Near a point of the higher
index locus, differentiating by weighted order less than $\sigma$
leaves weighted order at least $\epsilon$.  Multiplying by a
monomial cannot lower this order.  Restriction to $\Lambda$ is
factorwise and cannot lower it either.

Choose $c$ general linear combinations of these restored equations
so that their restrictions to $\Lambda$ have isolated intersections
at every indicated point.  This is possible because the common
base locus is isolated there: successive general combinations avoid
the finitely many associated components of positive dimension
through these finitely many points.  In completed local coordinates
$t_{\ell,1},\ldots,t_{\ell,c_\ell}$ at one such point, their
ideal is contained in the monomial ideal generated by monomials of
weighted order at least $\epsilon$.  The local intersection
multiplicity is therefore at least the number of monomials of
weight less than $\epsilon$.

In particular consider monomials for which every exponent in the
$\ell$th block is strictly less than $\epsilon D_\ell/(2c)$.
Their total weight is less than $\epsilon/2$, and their number is
at least
\begin{equation}\label{eq:cut-local-multiplicity}
 \left(\frac{\epsilon}{2c}\right)^c
       \prod_\ell D_\ell^{c_\ell}.
\end{equation}
Indeed there are $\lceil x\rceil\ge x$ nonnegative integers
strictly less than a positive real number $x$.

On the other hand, the sum of the isolated intersection
multiplicities of $c$ divisors of multidegree $D$ on $\Lambda$ is
at most
\begin{equation}\label{eq:cut-bezout}
 \left(\sum_\ell D_\ell H_\ell\right)^c\cdot\Lambda
 =\frac{c!}{\prod_\ell c_\ell!}\prod_\ell D_\ell^{c_\ell},
\end{equation}
where $H_\ell$ denotes the hyperplane class from the $\ell$th
factor.  The bound remains valid if other intersection components
have positive dimension.  Namely, perturb the divisors in their
basepoint-free complete linear systems to make the intersection
proper; at each original isolated complete intersection the local
length is conserved, and its contribution is part of the total
degree \eqref{eq:cut-bezout}.  There are
$q\prod_\ell\deg J_\ell$ of the isolated points counted above.
Combining their individual multiplicity bound with the total bound gives
\[
 q\left(\prod_\ell\deg J_\ell\right)
   \left(\frac{\epsilon}{2c}\right)^c
   \prod_\ell D_\ell^{c_\ell}
 \le \frac{c!}{\prod_\ell c_\ell!}
       \prod_\ell D_\ell^{c_\ell}.
\]
The powers of the multidegrees cancel, proving
\eqref{eq:conjugate-product-degree} independently of $D$.
For example, after harmless enlargement the constant is bounded by
$N!(2N/\epsilon)^N$.

Finally, the stabilizer of $J$ in the absolute Galois group of $k$
has index $q$.  Galois descent defines $J$ over its fixed field,
a degree-$q$ extension of $k$.  Each $J_\ell$ is the image of $J$
under a factor projection, so it is defined over the same field.
All factor degrees are positive integers; hence
\eqref{eq:conjugate-product-degree} bounds them separately as well.
\end{proof}

\subsection{The contradiction to minimality}

\begin{proposition}\label{prop:product-cut}
For the minimal products of Proposition~\ref{prop:minimal-products}
with $d_0>0$, the small sections of
Lemma~\ref{lemma:small-sections} cannot satisfy the positive index
conclusion of Proposition~\ref{prop:index} along an infinite
subsequence.
\end{proposition}

\begin{proof}
The finite projections and norms above give $G_i$ and $z_i$ with
\eqref{eq:cut-polynomial-index}.  Pass to constant dimensions of
the projective factors and let $N$ be their sum.  Apply the preceding
component construction with $\eta=\eta_1$.  Its $N+2$ levels and their
gap $\epsilon$ are fixed independently of $i$.  It gives a proper
component $J_i$ through $z_i$ common to two adjacent index loci;
the adjacent pair may depend on $i$.

For at least two factors, \eqref{eq:cut-ratios} eventually implies
\eqref{eq:cut-product-threshold}, so the product theorem gives
$J_i=\prod_\ell J_{i\ell}$.  With one factor this conclusion
requires no theorem.  Lemma~\ref{lemma:conjugate-product-degree}
then gives a common field of definition $k_i'$ with
$[k_i':k_i]$ bounded and all $\deg J_{i\ell}$ bounded.
The finitely many possible choices of adjacent levels do not
affect these bounds.

Because $J_i$ is proper, at least one factor $J_{i\ell}$ is
proper in $\mathbb P^{d_\ell}$.  After extraction fix this index
$\ell$.  A proper projective variety of bounded degree admits a
nonzero homogeneous equation of bounded degree over its field of
definition.  One can obtain it by projecting generally over that
field to a hypersurface in a projective space of one larger
dimension than the variety and pulling back its equation; the
image degree does not exceed the original degree.  Let $R_i$ be
such an equation for $J_{i\ell}$ over $k_i'$.

Pulling back $R_i$ through the linear projection $\pi_{i\ell}$
gives a bounded-degree hypersurface over $k_i'$ containing
$y_{i\ell}$.  It does not contain $Z_{i\ell}$, because that
projection is surjective onto $\mathbb P^{d_\ell}$ and $R_i$ is
nonzero.  The degrees $[k_i':\mathbb Q]$ are bounded, since the
degrees $[k_i:\mathbb Q]$ were bounded in choosing the minimal
products.  This is exactly the hypersurface excluded by
Proposition~\ref{prop:minimal-products}, a contradiction.
\end{proof}

\begin{proof}[Completion of the proof of Theorem~\ref{thm:sequence}]
Theorem~\ref{thm:all-high} handles the case with only unbounded
scales.  For a counterexample with a low factor, choose one of
minimal ambient dimension.  The projection and tangent arguments
apply in this minimal counterexample.  Proposition~\ref{prop:minimal-products}
then provides its auxiliary products.  If $d_0=0$,
Lemma~\ref{lemma:fixed-low-point} produces an all-high
counterexample, already impossible.  If $d_0>0$,
Proposition~\ref{prop:arithmetic-positive} and
Lemma~\ref{lemma:small-sections} provide the small sections;
Proposition~\ref{prop:index} gives their positive index; and
Proposition~\ref{prop:product-cut} contradicts minimality of the
auxiliary products.  Thus there is no counterexample with a low
factor either.
\end{proof}
